% Owner: guardsynth-coc / guardsynth-studio-sequential-coc
% Standalone formal specification; build with XeLaTeX twice.
% PDF and build outputs belong in the owner-scoped artifact directory.
\documentclass[11pt,a4paper]{article}
\usepackage[margin=24mm]{geometry}
\usepackage{fontspec}
\usepackage{kotex}
\setmainfont{Latin Modern Roman}
\setmainhangulfont[ItalicFont={Noto Sans CJK KR},ItalicFeatures={FakeSlant=0.15}]{Noto Sans CJK KR}
\setsanshangulfont[ItalicFont={Noto Sans CJK KR},ItalicFeatures={FakeSlant=0.15}]{Noto Sans CJK KR}
\usepackage{amsmath,amssymb,mathtools}
\usepackage{booktabs,tabularx,array}
\usepackage{xcolor}
\definecolor{accent}{RGB}{27,74,101}
\usepackage[colorlinks=true,linkcolor=accent,urlcolor=accent,citecolor=accent]{hyperref}
\hypersetup{pdftitle={GuardSynth Studio 검증 성질 정형 명세},pdfauthor={GuardSynth 연구 프로젝트}}
\setlength{\parindent}{0pt}
\setlength{\parskip}{0.4em}
\setlength{\emergencystretch}{1.5em}
\renewcommand{\arraystretch}{1.25}
\newcolumntype{Y}{>{\raggedright\arraybackslash}X}
\newcommand{\pred}[1]{\mathsf{#1}}
\newcommand{\status}[1]{\textnormal{\textsf{#1}}}
\newcommand{\E}{\mathcal{E}}
\newcommand{\G}{\mathbf{G}}
\newcommand{\F}{\mathbf{F}}
\newcommand{\sat}{\operatorname{Sat}}
\newcommand{\checksat}{\operatorname{Check}}
\newcommand{\supp}{\operatorname{Supp}}
\newcommand{\match}{\operatorname{Match}}
\newcommand{\obl}{\operatorname{Obl}}
\newcommand{\pid}[1]{\textbf{#1}}
\newcommand{\example}{\par\textbf{예와 판정.}\ }
\newcommand{\impl}{\par\textbf{구현 범위.}\ }
\renewcommand{\abstractname}{초록}
\renewcommand{\tablename}{표}
\renewcommand{\refname}{참고문헌}
\title{GuardSynth Studio 검증 성질 정형 명세\\[0.4em]
\large Safety, Consistency, Grounding, Coverage에 따른 B1--B8 및 A1--A4}
\author{GuardSynth 연구 프로젝트\\\small 문서 버전 v01 · 소유 프로젝트 \texttt{guardsynth-coc}}
\date{2026년 10월 4일}

\begin{document}
\maketitle

\begin{abstract}
녹화된 주행 장면으로부터 생성한 Chain of Causation(CoC)와 제약 후보에 대해
12개 검증 성질을 정형적으로 정의한다. 검증 단위는 바로 이전 장면과 현재 장면,
그리고 이전 장면에서 전달한 유한 상태 요약이다. 주분류는
Safety 5개, Consistency 3개, Grounding 3개, Coverage 1개이다.
각 성질의 입력, 적용 전제, 논리식과 판정 방법을 구분한다.
무기한 Liveness는 B6 및 B8의 확장으로 별도 기술하며,
조건의 만족 가능성과 비공허성은 추가 정의 검사로 분리한다.
본 문서는 목표 명세와 현재 구현 범위를 함께 설명한다.
검증 통과는 입력과 명세 안에서의 일치 또는 근거 충족을 의미한다.
\end{abstract}

\section{검증 대상과 분류}
CoC의 설명, 구조화된 관찰, 제약의 대상과 조건,
장면 간 유지 의무, 실제 관찰 행동을 검증한다.
Safety와 Liveness는 실행의 시간적 성질을 나타낸다.
Consistency, Grounding, Coverage는 검사의 대상과 목적을 나타낸다.
따라서 표~\ref{tab:classes}의 주분류는 배타적인 수학적 분해가 아니라
본 연구의 검증 체계를 설명하는 분류이다.
일관성 조건을 모든 시점에 요구하면 Safety 불변식으로도 표현할 수 있다.
시간적 구분은 문헌~\cite{lamport}에 따르며, 아래 분류와 운전 장면별 논리식은 본 연구의 설계이다.

\begin{table}[htbp]
\centering\small
\caption{12개 성질의 주분류와 관련 분류}\label{tab:classes}
\begin{tabularx}{\textwidth}{@{}l l Y Y@{}}
\toprule
ID & 주분류 & 검증 성질 & 관련 분류 또는 확장\\
\midrule
B5 & Safety & 조건 유지와 근거 있는 해제 & 상태 전이 일관성\\
B6 & Safety & 제약과 관찰 행동의 적합성 & 제약 일관성; 수행 의무의 Liveness 확장\\
B7 & Safety & 행동 선행 조건 충족 & Grounding\\
B8 & Safety & 기한 또는 결과 확인 시점의 후행 조건 & Grounding; 무기한 완료의 Liveness 확장\\
A3 & Safety & 개별 허용과 전체 허용의 구분 & 허용·금지 결합 일관성\\
\midrule
B1 & Consistency & 대상·형식·객체 대응 적합성 & 구조 검사\\
B3 & Consistency & 관찰·공간 관계 일관성 & 시간적 불변식으로 확장 가능\\
B4 & Consistency & 동일 ID·버전의 정의 유지 & 정의 변경 제한의 Safety\\
\midrule
B2 & Grounding & 근거의 시각·대상 적합성 & 판단 시점에 관한 Safety\\
A1 & Grounding & 설명에서 주장한 이유의 충족 & 주장과 근거 일관성\\
A4 & Grounding & 행동과 CAN 측정의 일치 & 측정값과 레이블 일관성\\
\midrule
A2 & Coverage & 검토된 필수 조건의 포함 & 명시적인 조건 목록에 상대적인 검사\\
\bottomrule
\end{tabularx}
\end{table}

검증은 추출된 입력의 품질에 의존한다.
같은 VLM이 서로 일치하는 잘못된 CoC와 제약을 만들면 내부 일관성만으로 이를 걸러낼 수 없다.
이 명세는 물리적 장면의 진실성, 사고가 발생하지 않은 원인의 인과적 식별,
행동의 유일한 정당성 또는 실제 주행 안전성에 대한 증명이 아니다.

\section{공통 모델과 판정 의미}
\subsection{두 장면과 상태 요약}
시점 인덱스를 $t$, 기록 시각을 $\tau_t$라 하고 $\tau_{t-1}<\tau_t$를 요구한다.
검증 창과 입력은 다음과 같다.
\begin{align}
 W_t&=(X_{t-1},X_t,\Sigma_{t-1}),\label{eq:window}\\
 X_t&=(CoC_t,Constraints_t,Evidence_t,\tau_t).
\end{align}
$\Sigma_{t-1}$은 조건 ID와 정의, 버전, 유지 의무,
이전 객체 대응, 직전 행동, 근거 참조와 관찰 공백 표시를 담는다.
전체 과거 관찰을 매번 SMT에 넣지 않는다.
첫 장면에서 이전 요약이 없으면 현재 장면부터 검사하고 이전 이력의 검증 범위를 표시한다.
유지 중인 의무를 요약 용량 때문에 조용히 삭제해서는 안 된다.

수식의 대상과 조건 집합은 유한하다.
본 명세의 $\forall$와 $\exists$는 해당 장면 창과 요약의 유한 집합에 대해 전개할 수 있다.
별도 양화가 없는 요구식의 자유변수는 해당 검증 범위의 대상에 대해 보편적으로 적용한다.
과거와 현재의 원자는 시점 인덱스를 구분한다.
과거의 신호색이나 차량 위치를 현재 사실로 복사하지 않는다.

\subsection{대상과 행동 및 제약}
행동은 다음 튜플이다.
\begin{equation}
 \begin{aligned}
 a&=(code,actor,target,movement,I_a),\\
 code&\in\{\pred{START},\pred{STOP},\pred{CONTINUE},\pred{ENTER},\pred{HOLD}\}.
 \end{aligned}
\end{equation}
$I_a$는 관찰 또는 의무의 시간 범위이다.
$\pred{HOLD}$는 이미 정지한 상태의 유지,
$\pred{CONTINUE}$는 주행 지속, $\pred{STOP}$은 정지 동작을 뜻한다.
정지 동작을 시작한 시점에 이미 속도가 0이어야 한다고 요구하지 않는다.
$\pred{ENTER}$와 $\pred{CONTINUE}$는 동시에 기록될 수 있다.

제약 정의의 의미 필드는 다음과 같다.
\begin{align}
 d={}&(id,version,class,actor,T_d,target,movement,\notag\\[-0.3em]
      &\quad A_d,effect,G_{on},G_{off},Pre,Post,scope,source),\label{eq:definition}\\
 effect&\in\{\pred{FORBID},\pred{REQUIRE},\pred{ALLOW}\}.
\end{align}
$T_d$는 조건의 대상 집합, $A_d$는 행동 종류 집합이다.
$\match(d,a,t)$는 행위자, 행동 종류, 구체 대상, 진행 방향 및 적용 시간 범위의 대응이 확인되었음을 뜻한다.
미지정 대상을 임의로 모든 대상으로 확장하지 않는다.
현재 구현에 없는 버전·시간 구간·완료 기한 등은 설계상 필드이며,
그 필드가 필요한 성질의 검사는 입력이 갖추어질 때 수행한다.

각 원자 조건 $p$와 요구 진릿값 $b$의 리터럴을
\begin{equation}
 \ell(p,b)=\begin{cases}p,&b=\pred{TRUE},\\ \neg p,&b=\pred{FALSE}\end{cases}
\end{equation}
로 정의한다. 활성, 해제, 선행 조건의 결합은
조건항 $g=(p,b)$에 대한 약기 $\ell(g)=\ell(p,b)$를 사용하여
\begin{equation}
 G_{on}=\bigwedge_{g\in on(d)}\ell(g),\qquad
 G_{off}=\bigvee_{g\in off(d)}\ell(g),\qquad
 Pre=\bigwedge_{g\in pre(d)}\ell(g)
\end{equation}
이다. 빈 활성 목록은 입력 부족으로 처리한다.
빈 해제 목록은 영구 조건인지 정의 누락인지 검토하고,
자동으로 잘못된 제약이나 완료된 제약으로 판정하지 않는다.

\subsection{관찰과 근거 모델}
관찰값은 $\{\pred{TRUE},\pred{FALSE},\pred{UNKNOWN},\pred{CONFLICT}\}$이다.
$\pred{UNKNOWN}$은 거짓이 아니라 결정되지 않은 값이다.
인정된 관찰값만 다음 사실식에 넣는다.
\begin{equation}
 F_t=\bigwedge_{p:v_t(p)=\pred{TRUE}}p_t
       \land\bigwedge_{p:v_t(p)=\pred{FALSE}}\neg p_t.
\end{equation}
미확인 원자는 참·거짓으로 고정하지 않는다.
같은 대상·시점·좌표 기준의 사실이 충돌하면 입력 충돌을 먼저 보고한다.

근거 모델 $\E_t$는 시점별 관찰 사실, 검토된 관계 의미,
객체 대응과 시간 정합성을 결합한다.
검사하려는 설명의 주장이나 기대 결과를 그 자체로 근거에 추가하지 않는다.
규범적 제약 $K_t$, 상태 갱신 $T_t$와 관찰 행동 $O_t$는 근거 사실과 구분한다.
미기록 행동은 관찰 완전성이나 명시적인 부정 관찰이 없으면 $O_t$에서 거짓으로 고정하지 않는다.

\subsection{함의 검사와 양립 가능성 검사}
$\checksat(\psi)\in\{\status{SAT},\status{UNSAT},\status{UNKNOWN}\}$를 SMT 결과라 하고,
$\sat(\psi)$는 결과가 $\status{SAT}$임을 뜻한다.
근거가 일관된 경우의 주장 판정은 다음과 같다.
\begin{equation}
 \operatorname{Verdict}_t(\phi)=
 \begin{cases}
 \status{SUPPORTED},&\checksat(\E_t\land\neg\phi)=\status{UNSAT},\\
 \status{CONTRADICTED},&\checksat(\E_t\land\phi)=\status{UNSAT},\\
 \status{UNCONFIRMED},&\checksat(\E_t\land\phi)=\status{SAT}\ \text{및}\\
 &\checksat(\E_t\land\neg\phi)=\status{SAT},\\
 \status{SOLVER\_UNKNOWN},&\text{필요한 쿼리가 미결정인 경우}.
 \end{cases}\label{eq:verdict}
\end{equation}
먼저 $\checksat(\E_t)=\status{UNSAT}$이면
$\status{INPUT\_CONFLICT}$로 처리하여 모순된 전제로 다른 주장을 증명하지 않는다.
이 부정식을 이용한 함의 판정은 문헌~\cite{z3guide}의 방식에 따른다.

$\supp_t(\phi)$는 일관된 $\E_t$에 대해 $\E_t\models\phi$가 확인되었음을 뜻하는
\emph{검사 수준의 술어}이다. 아래 수식에서 $\supp$, $\sat$, 참조 검사 등은
상태 원자와 구분되는 메타 검사이며 각각 함의 쿼리, SAT 쿼리, 구조 검사로 구현한다.
근거와 주장 양쪽이 가능하면 반례 후보가 존재해도 확정 오류로 세지 않는다.

입력 필드나 연결이 없으면 $\status{NOT\_CHECKED}$,
대상이 아니면 $\status{NOT\_APPLICABLE}$로 표시한다.
조건부 성질의 적용 전제가 미확인이면 적용성을 보류한다.
전제가 성립하지 않았다는 이유로 그 성질을 의미 있게 검증했다고 표시하지 않는다.

\clearpage
\section{Safety에 속하는 성질}
Safety는 금지된 기록이나 전이를 허용하지 않는 요구이다.
일관되고 완전한 모델의 실행 $\sigma$에서는 대표적으로
$\sigma\models\G\neg Bad$로 표현한다.
아래 식은 해당 모델의 유한한 장면 창에서 검사할 국소 요구이다.
부분 관찰에서는 미확인 입력을 보류하며 전체 실행의 안전성을 추정하지 않는다.

\subsection{B5 조건 유지와 근거 있는 해제}\label{prop:b5}
활성 및 해제 확인을 다음처럼 정의한다.
\begin{equation}
 c^{on}_{d,t}:=\supp_t(G_{on}(d,t)),\qquad
 c^{off}_{d,t}:=\supp_t(G_{off}(d,t)).
\end{equation}
동시 활성·해제 또는 관련 근거 충돌을 $Conflict(d,t)$라 한다.
충돌 없는 갱신은 다음을 만족해야 한다.
\begin{align}
 &\neg Conflict(d,t)\land c^{off}_{d,t}
   \Rightarrow Released(d,t)\land\neg\obl(d,t),\label{eq:b5release}\\
 &\neg Conflict(d,t)\land\neg c^{off}_{d,t}
   \land\bigl(\obl(d,t-1)\lor c^{on}_{d,t}\bigr)
   \Rightarrow\obl(d,t),\label{eq:b5carry}\\
 &\neg\bigl(c^{on}_{d,t}\land c^{off}_{d,t}\bigr).
\end{align}
이전 의무가 있는 상태에서 충돌이 생기면 의무를 검토 대기 상태로 보존한다.
$\neg c^{off}_{d,t}$는 해제를 확인하지 못했다는 뜻이며,
$\supp_t(\neg G_{off})$와 같지 않다.
의무 유지는 물리적 장애물의 현재 존재를 관찰했다는 주장이 아니다.

\example 이전 장면의 보행자 진입 보류 조건이 현재 목록에서 빠져도,
해제 근거가 확인되지 않았다면 이전 요약에서 유지한다.
처음부터 활성 여부를 알지 못한 조건을 활성 의무로 만들어내지 않는다.
\impl 유한 원장과 유지·해제 상태 갱신은 구현되어 있다.
해제 미확인 표시와 논리적 모순 판정을 구분한다.

\subsection{B6 제약과 관찰 행동의 적합성}\label{prop:b6}
행동 원자를 $Do(a,t)$, 정해진 구간 안의 수행을 $Done(a,I_d)$라 한다.
\begin{align}
 &\obl(d,t)\land\match(d,a,t)\land effect(d)=\pred{FORBID}
   \Rightarrow\neg Do(a,t),\label{eq:b6forbid}\\
 &\obl(d,t)\land\match(d,a,t)\land effect(d)=\pred{REQUIRE}\notag\\[-0.3em]
 &\hspace{3em}\land Due(d,t)\Rightarrow Done(a,I_d),\label{eq:b6require}\\
 &effect(d)=\pred{ALLOW}\ \not\Rightarrow Do(a,t).
\end{align}
$Due(d,t)$는 수행 의무의 판정 시점에 도달했음을 뜻한다.
해당 시점 의무라면 현재 시점, 구간 의무라면 명시한 기한을 사용한다.
두 단계의 양립 가능성 쿼리는
\begin{align}
 Q_{rules}&:=\sat(\E_t\land T_t\land K_t),\\
 Q_{record}&:=\sat(\E_t\land T_t\land K_t\land O_t)
\end{align}
이다. 일관된 근거 아래 $Q_{rules}$가 거짓이면 제약 간 충돌이다.
$Q_{rules}$는 참인데 $Q_{record}$가 거짓이면 관찰 행동과 제약의 충돌이다.
SAT는 양립 가능성의 결과이며, 모든 미확인 조건의 충족 증명은 아니다.

\example 같은 영역의 진입을 요구하면서 금지하면 제약 간 충돌이다.
녹색 신호의 허용 후보와 시야 부족의 금지가 공존한다는 사실만으로
두 제약을 모순으로 판정하지 않는다.
\impl 활성 요구·금지와 기록 행동의 대조는 구현되어 있다.
현재 행동 쿼리는 미기록 행동을 거짓으로 설정하므로,
관찰 누락과 미수행을 구분하는 본 명세와 차이가 있다.
시간 구간을 갖는 수행 의무는 추가 모델링이 필요하다.

\subsection{B7 행동 선행 조건의 충족}\label{prop:b7}
$\mathcal{P}_{req}(a,t)$는 해당 행동에 필수로 적용되는 선행 조건의 집합이다.
활성 의무, 검토된 행동 조건, 또는 명시적으로 선택한 허용 근거에서 결정한다.
선택되지 않은 다른 허용 후보의 조건을 모든 행동에 강제하지 않는다.
\begin{equation}
 Obs(a,t)\land PreScopeResolved(a,t)
 \Rightarrow\bigwedge_{p\in\mathcal{P}_{req}(a,t)}\supp_{pre(a)}(p).
 \label{eq:b7}
\end{equation}
$\E_{pre(a)}$는 행동 판단 전에 이용 가능하며 해당 대상·시각에 적합한 근거 모델이다.
필수 조건이 거짓으로 확인되면 불일치, 결정되지 않으면 미확인이다.
필수 목록의 연결이 없으면 검사 불가로 남긴다.
목록의 충분성은 A2에서 별도로 검사한다.

\example 적용되는 진입 선행 조건이 $PATH\_CLEAR(z)$인데
해당 경로의 비어 있지 않음이 확인되면 불일치이다.
경로가 영상에 가려졌다면 미확인이다.
\impl 활성 제약에 입력된 선행 조건과 기록 행동의 대조를 제공한다.
의사결정 이전 근거의 상세 시간 구간과 필수 조건 선택 의미는 부분적으로만 모델링되어 있다.

\subsection{B8 행동 후 기대 결과의 적합성}\label{prop:b8}
행동과 후행 조건이 연결되고, 결과를 평가할 구간이 종료되었을 때 검사한다.
\begin{align}
 &Obs(a,t-1)\land LinkedPost(d,a)\land SameTarget(d,a,t)\notag\\[-0.3em]
 &\hspace{2em}\land OutcomeWindowClosed(d,a,t)
 \Rightarrow\supp_t(Post(d,a,t)).\label{eq:b8}
\end{align}
명확한 반대 결과는 불일치, 결과 관찰 부족은 미확인이다.
허용되었다는 사실만으로 수행이나 완료를 요구하지 않는다.
제약의 해제와 행동의 완료는 서로 다른 상태이다.

\example ``정지 동작을 시작했다''는 설명만으로 다음 썸네일에서
속도 0을 요구하지 않는다. 정지 완료 시점이나 명시한 완료 기한을 사용한다.
무기한 완료 요구는 절~\ref{sec:liveness}의 Liveness 확장이다.
\impl 이전 행동의 기대 결과와 바로 다음 장면을 대조한다.
명시적인 결과 확인 구간·완료 기한 입력은 없으므로,
현재 검사는 다음 장면에서 평가 가능한 결과에 한정해야 한다.

\subsection{A3 개별 허용과 전체 허용의 구분}\label{prop:a3}
현재의 개별 허용 후보가 충족되는 식을 다음과 같이 정의한다.
\begin{align}
 AllowReady(d,a,t):={}&\match(d,a,t)\land CurrentScope(d,t)\notag\\[-0.3em]
 &\land Active(d,t)\land G_{on}(d,t)\land\neg G_{off}(d,t)\land Pre(d,a,t).
\end{align}
검토 대상 제약 집합 안에서의 전체 허용 근거는
\begin{align}
 Permit(a,t):={}&
 \left(\bigvee_{d:effect(d)=\pred{ALLOW}}AllowReady(d,a,t)\right)\notag\\[-0.3em]
 &\land\neg\left(
 \bigvee_{\substack{f:effect(f)=\pred{FORBID}\\\match(f,a,t)}}\obl(f,t)
 \right).\label{eq:a3}
\end{align}
이다. 이 식에 식~\eqref{eq:verdict}의 근거 판정을 적용한다.
현재 허용 전제 또는 금지 여부가 미확인이면 전체 허용을 확정하지 않는다.
유지 중인 금지가 확인되면 개별 허용 후보가 충족되어도 전체 허용은 차단된다.
\begin{equation}
 Permit(a,t)\not\Rightarrow Do(a,t),\qquad
 Permit(a,t)\not\Rightarrow\pred{REQUIRE}(a,t).
\end{equation}
식~\eqref{eq:a3}은 예외가 없는 프로필을 기준으로 한다.
예외와 우선 규칙은 검토된 정의로만 적용한다.
제약 집합 밖의 누락 가능성은 A2에서 다루며, 이 판정을 실제 주행 허가로 사용하지 않는다.

\example 녹색 신호가 자차에 적용되어도 해당 진입의 시야 부족 금지가 유지되면
전체 허용은 차단된다. 실제 행동이 정지 유지라면 그 사실은 행동 위반이 아니다.
\impl 전체 허용의 추가 판정은 오프라인 실험 단계이다.
현재 생성·검증 경로에 완전히 통합된 성질로 해석하지 않는다.

\section{Consistency에 속하는 성질}
\subsection{B1 대상과 관계의 형식 적합성}\label{prop:b1}
$\mathcal{P}_t$와 $\mathcal{A}_t$를 사용 원자와 행동의 집합으로 두고,
미지정 참조를 $\bot$로 표시한다.
\begin{align}
 \forall p\in\mathcal{P}_t:\quad
 &Bound(subj(p))\land\bigl(obj(p)\ne\bot\Rightarrow Bound(obj(p))\bigr)\notag\\[-0.3em]
 &\land ArityKindCompatible(p)\land IdentityResolved(p),\label{eq:b1atom}\\
 \forall a\in\mathcal{A}_t:\quad
 &ActorBound(a)\land TargetCompatible(a).\label{eq:b1action}
\end{align}
$ArityKindCompatible$은 단항·이항 관계 및 대상 종류의 적합성을 뜻한다.
$IdentityResolved$는 장면 간 같은 객체로 연결하는 대응을 확인했다는 뜻이며,
불확실한 대응을 임의로 확정하지 않는다.

\example $SIGNAL\_GREEN(s)$는 신호기,
$PATH\_CLEAR(z)$는 차로나 영역에 관한 원자이다.
대상이 자차만으로 지정된 경로 여유 조건으로 특정 교차로 진입을 정당화할 수 없다.
\impl 종류·참조 검사와 불확실한 객체 대응 처리를 제공한다.
객체 동일성의 물리적 정확성은 별도 검토가 필요하다.

\subsection{B3 관찰과 공간 관계의 일관성}\label{prop:b3}
같은 시점·좌표 기준의 관계 의미를 $R_t$라 하면
\begin{equation}
 \sat(F_t\land R_t)\label{eq:b3}
\end{equation}
를 요구한다. 좌우 관계의 대표 공리는
\begin{align}
 LEFT(x,y,t)&\Leftrightarrow RIGHT(y,x,t),\\
 \neg LEFT(x,x,t),\qquad
 \neg\bigl(LEFT(x,y,t)&\land LEFT(y,x,t)\bigr),\\
 LEFT(x,y,t)\land LEFT(y,z,t)&\Rightarrow LEFT(x,z,t)
\end{align}
이다. 전후 관계도 엄격한 순서로 모델링한다.
유한 대상에 대해 정수 위치 $u_t(x)$를 사용하면
\begin{equation}
 AHEAD(x,y,t)\Leftrightarrow u_t(x)>u_t(y)
\end{equation}
로 표현할 수 있다. 다른 시점이나 다른 좌표 기준의 관계를 한 원자로 합치지 않는다.

\example 실제 과거 생성본에 등장한 $AHEAD(z,z)=\pred{TRUE}$는
$u(z)>u(z)$를 요구하여 불가능하다.
원래 의미가 교차로가 자차 앞에 있다는 것이라면 후보 수정은 $AHEAD(z,ego)$이다.
$AHEAD(ego,z)$는 자차가 교차로보다 앞에 있다는 다른 의미이다.
후보 수정의 SAT가 영상의 정확성을 확정하지는 않는다.
\impl 현재 시점의 순서 축으로 역관계, 반대 방향과 순환 모순을 검출한다.
좌표 기준을 명시하는 입력 필드는 아직 없다.

\subsection{B4 제약 정의의 동일성}\label{prop:b4}
조건의 ID와 버전이 같으면 구조화된 의미 정의를 유지한다.
\begin{align}
 &id(d_{t-1})=id(d_t)\land version(d_{t-1})=version(d_t)\notag\\[-0.3em]
 &\hspace{3em}\Rightarrow Def(d_{t-1})=Def(d_t),\label{eq:b4}\\
 &DefChanged(d)\Rightarrow NewVersionOrID(d)\land ChangeRecorded(d).
\end{align}
$Def$에는 식~\eqref{eq:definition}의 대상, 효과, 행동,
활성·해제·선행·후행 조건과 적용 범위 등이 포함된다.
정의와 현재 사실을 분리하여, 신호색이나 차량 위치의 갱신을 정의 변경으로 처리하지 않는다.

\example 같은 조건 ID로 이전에는 공사 구역 진입을 금지하고,
현재는 다른 차로의 주행을 요구하도록 바꾸면 정의 충돌이다.
\impl 현재 구현은 같은 ID의 전체 정의를 구조적으로 비교한다.
명시적인 버전 전환과 의미 필드 중심의 변경 관리는 추가 설계 사항이다.

\section{Grounding에 속하는 성질}
\subsection{B2 근거의 시간과 적용 대상 적합성}\label{prop:b2}
사실로 사용하는 주장에는 허용된 근거 연결이 있어야 한다.
\begin{align}
 Accepted(p_t)\Rightarrow\exists e:\quad
 &Linked(e,p_t)\land SupportsKind(e,p_t)\notag\\[-0.3em]
 &\land ScopeMatch(e,p_t)\land TimeAdmissible(e,p_t).\label{eq:b2fact}
\end{align}
$SupportsKind$는 해당 종류의 근거로 원자를 관찰할 수 있음을 뜻한다.
CAN만으로 신호색이나 영역 진입을 관찰했다고 처리하지 않는다.
현재 상태에 대한 근거의 유효 기간은 명시적으로 설정해야 한다.
판단 이유의 시간 요구는
\begin{equation}
 UsedAsReason(e,a)\Rightarrow end(e)\le\tau_{decision}(a)\label{eq:b2time}
\end{equation}
이다. 시간 범위가 겹치거나 판단 시각이 없으면 이를 확정하지 않는다.
후행 결과의 근거는 행동 뒤에 올 수 있지만 당시 판단 이유로 재사용하지 않는다.

\example 다음 장면에서 교차 차량이 사라졌다는 사실은
이전 진입 판단 당시 경로가 비어 있었다는 근거가 될 수 없다.
\impl 현재 근거 ID와 원자별 근거 종류를 검사한다.
판단 시각·근거 유효 기간·문장별 시간 연결은 부분 구현이다.
근거 ID가 있다는 사실이 VLM의 영상 해석을 인증하지 않는다.

\subsection{A1 설명에서 주장한 이유의 근거 충족}\label{prop:a1}
$J$는 행동의 이유로 명시한 유한한 조건 집합이며
$\bigwedge J$는 그 이유의 결합이다.
\begin{equation}
 ExplainReason(a,J)\Rightarrow
 Bound(J,a)\land Relevant(J,a)\land\supp_{pre(a)}\!\left(\bigwedge J\right).
 \label{eq:a1}
\end{equation}
$Bound$는 행동과 구체 원자의 연결,
$Relevant$는 이유가 해당 행동 판단에 관련된다는 검토된 대응이다.
빈 이유 연결을 공허한 증명으로 처리하지 않는다.
조건을 참으로 확인해도 그 조건이 실제 행동을 발생시킨 유일한 원인이라는 결론은 도출하지 않는다.

\example ``시야가 충분하고 진입 경로가 비어 있어 진입했다''는 설명은
시야 대상 $v$와 진입 영역 $z$에 대한 다음 조건에 연결되어야 한다.
\[
 J=\{\neg\pred{VISIBILITY\_INSUFFICIENT}(v),\ \pred{PATH\_CLEAR}(z)\}.
\]
이유와 반대인 관찰이 확인되면 설명의 불일치이다.
\impl 현재 CoC 형식에는 이유--행동--원자 조건의 명시적 연결 필드가 없다.
전체 설명에 대한 자동 검사는 이 연결이 추가되어야 가능하다.
오프라인 명시적 주장 시험과 실제 생성 설명의 오류 검출 효과를 구분한다.

\subsection{A4 관찰 행동과 CAN 측정의 일치}\label{prop:a4}
측정 $q$의 시간, 단위, 품질과 연결을 먼저 검사한다.
\begin{align}
 EligibleCAN(q,t):={}&Linked(q,t)\land UnitsValid(q)\land QualityAccepted(q)\notag\\[-0.3em]
 &\land LatestPastSample(q,t)\land 0\le\tau_t-\tau_q\le\delta
 \land speed(q)\ge0.
\end{align}
정지 경계 $\epsilon_s$와 이동 경계 $\epsilon_m$은
$0\le\epsilon_s<\epsilon_m$을 만족한다.
\begin{equation}
 DecidableMotion(q):=speed(q)\le\epsilon_s\lor speed(q)\ge\epsilon_m.
\end{equation}
$EligibleCAN\land DecidableMotion$이 확인된 측정에 한해
\begin{align}
 Obs(\pred{HOLD}(ego),t)&\Rightarrow speed(q)\le\epsilon_s,\label{eq:a4hold}\\
 Obs(\pred{CONTINUE}(ego),t)&\Rightarrow speed(q)\ge\epsilon_m.\label{eq:a4continue}
\end{align}
두 경계 사이의 측정값, 오래되거나 연결되지 않은 측정은 판정을 보류한다.
실험 프로토콜은 $\epsilon_s=0.5\,\mathrm{km/h}$,
$\epsilon_m=2.0\,\mathrm{km/h}$, $\delta=0.6\,\mathrm{s}$를 사용했다.
이 값은 보편적인 차량 기준이 아니라 해당 실험의 선택이다.

\example 유효하게 연결한 속도 $8.3\,\mathrm{km/h}$와
$\pred{HOLD}$ 레이블은 불일치이다.
이는 연결된 측정과 레이블의 대조이며 영상 시각까지 차량 상태가 불변이었다는 증명이 아니다.
\impl CAN 수치와 레이블의 추가 검사는 오프라인 실험 단계이다.
$\pred{START}$ 및 $\pred{STOP}$에는 충분한 전환 구간 측정이 필요하고,
$\pred{ENTER}$는 속도만으로 확인할 수 없다.

\section{Coverage에 속하는 성질}
\subsection{A2 장면과 행동에 필요한 조건의 포함}\label{prop:a2}
$Required(ctx,a,v)$는 생성 출력과 독립적으로 검토한
장면·행동별 필수 조건 목록과 그 버전이다.
$Covered$는 같은 대상, 요구 진릿값, 진행 방향과 적용 시간 범위에 맞게 정의된 조건이다.
\begin{align}
 Required(ctx_t,a,v)&\subseteq
 Covered(Definitions_t\cup\Sigma_{t-1},a),\label{eq:a2}\\
 Missing(a)&:=Required(ctx_t,a,v)-Covered(Definitions_t\cup\Sigma_{t-1},a),\\
 Missing(a)&=\varnothing.
\end{align}
참조와 종류가 맞지 않거나, 해당 행동 이전에 필요하지만 후행·해제 조건에만 등장하는 조건은 포함된 것으로 세지 않는다.
이 검사는 조건의 구조적 포함을 검사한다.
조건이 실제로 참인지는 B7 또는 A1의 근거 검사에서 따로 판정한다.
검토된 필수 목록이 없으면 Coverage를 검사할 수 없다.

\example 구체 교차로의 진입 조건 목록이 해당 경로의 여유를 요구한다면,
녹색 신호 조건만 있는 정의에는 $PATH\_CLEAR(target(a))$가 누락되어 있다.
신호의 자차 적용성이 필요한 목록이라면 그 연결도 별도로 요구한다.
이러한 목록은 검토한 범위에 상대적이며 모든 국가·장면의 교통 규칙을 완전히 열거했다는 뜻은 아니다.
\impl 오프라인 실험에서는 진입·출발 대상의 경로 조건과
녹색 신호의 적용 대상 연결을 검사했다.
해당 실험 기준과 전체 주행 정당화의 완전성을 구분한다.

\section{Liveness의 조건부 확장}\label{sec:liveness}
위 12개 국소 성질에는 무기한 진행 보장이 독립적으로 포함되어 있지 않다.
B6의 수행 의무와 B8의 완료 요구를 확장하려면
실행 모델과 환경 가정 $\mathcal{H}$를 먼저 정해야 한다.
같은 요청 또는 행동 에피소드 $r$에 대해 다음을 사용할 수 있다.
\begin{align}
 Spec\land\mathcal{H}&\models
 \G\bigl(Required(r)\Rightarrow\F Done(r)\bigr),\label{eq:liveb6}\\
 Spec\land\mathcal{H}&\models
 \G\bigl(Started(r)\Rightarrow\F Completed(r)\bigr).\label{eq:liveb8}
\end{align}
요청의 취소나 적용 범위 종료를 허용하는 경우에는
그 의미를 완료 조건 또는 별도의 종료 조건으로 명시해야 한다.
같은 행동 종류의 다른 에피소드 완료를 $r$의 완료로 사용하지 않는다.
허용 후보만으로 $Required$를 생성하지 않는다.

적색 신호나 장애물 차단이 녹화 종료까지 계속되는 실행에
무조건적인 출발·진입·해제를 요구하면 부적절하다.
공정성은 계속 실행 가능하게 유지되는 동작이 무한히 무시되지 않도록 하는 가정이다~\cite{lamport}.
환경의 신호 변경이나 장애물 소멸을 공정성만으로 만들어내지 않는다.
실행 가능 조건이 계속 유지될 것인지도 환경과 전이 모델에서 확인해야 한다.

기한 $\Delta$를 명시한 요구는
\begin{equation}
 \G\bigl(Triggered(r)\Rightarrow\F_{\le\Delta}Completed(r)\bigr)
 \label{eq:boundedresponse}
\end{equation}
처럼 표현할 수 있다.
기한을 넘긴 시점에 위반을 확인할 수 있으므로 유한한 반례를 갖는 Safety 요구이다.
두 장면과 시작 시각·완료 상태의 요약으로 이 기한을 추적할 수 있으나,
관찰이 불완전하면 미완료를 확정하지 않고 보류한다.
무기한 Liveness의 일반적인 보장을 두 장면만으로 증명하지 않는다.

\section{추가 정의 검사와 검증 운영}
\subsection{실행 가능성과 비공허성}
이 절의 정의 검사는 12개 성질 외에 논의한 추가 검사이다.
현재 관찰 사실을 고정하지 않고 해당 정의를 만족하는 가상 상태가 있는지 검사한다.
\begin{align}
 \sat(Domain_d\land G_{on}(d)),\qquad
 \sat(Domain_d\land G_{off}(d)),\label{eq:feasibleguards}\\
 \sat(Domain_d\land G_{on}(d)\land Pre(d)),\label{eq:feasibleexecution}\\
 \sat(Domain_d\land G_{on}(d)\land\neg G_{off}(d)).\label{eq:effectiveactivation}
\end{align}
식~\eqref{eq:feasibleguards}의 해제 가능성은 해제가 가능해야 한다고 명세한 조건에 적용한다.
영구 조건에도 해제를 강제하지 않는다.
식~\eqref{eq:feasibleexecution}은 행동에 적용되는 선행 조건과 활성 조건의 양립 가능성이다.
식~\eqref{eq:effectiveactivation}은 동시에 해제되지 않고 활성화할 수 있는지 보는 진단이다.
활성·해제 우선 의미가 있으면 그 의미를 반영해야 한다.

$G_{on}=p\land\neg p$는 불가능하다.
반면 해제 목록의 $p$와 $\neg p$는 $G_{off}=p\lor\neg p$가 되어 항상 참이다.
이는 해제 불가능성과 다른 문제이다.
도달 가능한 전이 경로가 있는지, 모든 허용 실행에서 결국 해제되는지는
상태식의 SAT만으로 알 수 없다.
활성화가 불가능하면 $\G(Active\Rightarrow\F Released)$가 공허하게 통과할 수 있다.

실제 과거 생성 사례의 자기 순서 선행 조건은
식~\eqref{eq:feasibleexecution}을 불가능하게 만들었다.
활성·해제 조건 자체의 불가능성을 발견한 사례와 구분한다.
원본과 대상 수정안의 대조, 오류 주입 및 합성 반복 경로의 결과는
프로젝트의 실행 가능성 실험 기록에 보존되어 있다.

\subsection{재검증과 사람의 수정}
검증 결과에는 입력, 이전 상태 요약, 규칙 출처 목록,
프로필과 컴파일러 버전의 의존성을 연결한다.
이전 장면의 입력이 생성되거나 수정되면 해당 장면과 바로 다음 장면의 결과를 무효화한다.
상태 요약 변화가 이어지면 다음 장면쌍을 순차적으로 재검사한다.
각 SMT 문제의 범위는 여전히 식~\eqref{eq:window}의 두 장면이다.

불일치 결과에는 성질 ID, 장면, 대상, 실제 비교값,
입력 경로와 관찰 근거를 붙인다.
양립 불가능한 제약에는 충돌에 관여한 assertion의 unsat core를 연결한다~\cite{z3guide}.
이는 항상 최소 원인 집합인 것은 아니다.
가능한 반례 모델은 확정된 실제 오류와 구분한다.
사람이 수정한 뒤 의존하는 장면쌍을 다시 검사한다.
규칙 원문 검토 상태와 논리 판정을 별도 관리한다.

\subsection{구현과 효과에 관한 해석}
B1, B3, B4, B5 및 입력된 선행 조건의 B7은 현재 검증 경로에서 다룬다.
B2, B6, B8에는 시간·관찰 완전성 관련 차이가 있으며 본문의 구현 범위에 명시했다.
A1은 이유 연결 형식이 필요하고, A2--A4의 추가 판정은 오프라인 실험 단계이다.
본 문서의 정의가 모든 구현의 완성이나 실제 생성 오류 검출 효과를 의미하지 않는다.

기존의 추가 성질 1--4 평가에서 실제 원본의 새 확정 모순은 0건이었다.
후속 실행 가능성 검사에서 찾은 자기 순서 선행 조건은 추가로 살펴본 과거 생성본의 사례이다.
현재 프로필에서 이미 통과한 입력의 새로운 탈락으로 주장하지 않는다.
주입 오류의 검출 능력, 실제 생성 오류의 검출률,
정상 사례의 오탐률을 각각 구분하여 평가해야 한다.

\begin{thebibliography}{9}
\bibitem{lamport}
Leslie Lamport, \emph{Specifying Systems}, preliminary draft, 4 March 2002,
Chapter 8, ``Liveness and Fairness.''
\url{https://lamport.azurewebsites.net/tla/book-02-03-04.pdf}.
\bibitem{z3guide}
Z3 Guide, ``Propositional Logic,'' 2026년 10월 4일 확인.
\url{https://microsoft.github.io/z3guide/docs/logic/propositional-logic/}.
\end{thebibliography}
\end{document}
